\section*{Conclusion générale}
\addcontentsline{toc}{section}{Conclusion générale}
\label{sec:conclusion}

L'analyse conduite dans ce rapport permet de tirer plusieurs enseignements majeurs sur la comparaison entre assurance-vie et compte-titres ordinaire (CTO).
\begin{itemize}
    \item En ligne directe, qui constitue la configuration de loin la plus fréquente, l'assurance-vie ne parvient pas à compenser l'avantage structurel du CTO. La neutralisation des plus-values latentes au décès, conjuguée à l'absence de frais d'enveloppe récurrents, confère au CTO une efficacité fiscale supérieure dans l'immense majorité des cas.
    \item En ligne indirecte, l'assurance-vie conserve un avantage lié à son abattement spécifique de 152 500 euros et à des taux réduits. Bien que si l'espérance de vie de l'investisseur est très grande ou que les frais liées à l'assurance vie sont très importants, alors cet avantage pourrait être remis en cause.
    \item Les retraits programmés, même optimisés après huit ans, n'inversent pas cette hiérarchie : l'impact positif de l'abattement annuel demeure marginal face au poids des frais et des prélèvements sociaux.
    \item Le facteur décisif reste le niveau des frais. Les simulations montrent que, pour qu'une assurance-vie soit réellement compétitive face au CTO, ses frais devraient descendre sous 0,25\% par an de frais d'enveloppe — un niveau largement irréaliste dans le marché actuel, où les frais moyens dépassent nettement ce seuil \cite{rapport_observatoire_epargne_2025}.
\end{itemize}

En conclusion, il apparaît que, dans la quasi-totalité des situations rencontrées en pratique — transmissions en ligne directe et contrats assortis de frais élevés — l'assurance-vie n'est pas l'enveloppe la plus favorable pour les épargnants. La popularité persistante de ce produit semble davantage s'expliquer par la rente qu'il génère pour l'industrie financière, grâce à des frais captés de manière récurrente, que par un véritable avantage économique pour les souscripteurs. Ce succès repose ainsi sur une méconnaissance entretenue du cadre fiscal, au détriment de l'intérêt réel des ménages.
